\documentclass[letterpaper,11pt]{article}

\usepackage{iftex}
\usepackage{fontspec}
\usepackage[empty]{fullpage}
\usepackage[usenames,dvipsnames]{color}
\usepackage[hidelinks]{hyperref}
\usepackage{fancyhdr}
\usepackage[english]{babel}
\usepackage{setspace}

\setmainfont[
  Path=fonts/,
  UprightFont=Geist-Regular.ttf,
  BoldFont=Geist-SemiBold.ttf,
  ItalicFont=Geist-Regular.ttf,
  BoldItalicFont=Geist-SemiBold.ttf,
  ItalicFeatures={FakeSlant=0.2},
  BoldItalicFeatures={FakeSlant=0.2}
]{Geist}

\pagestyle{fancy}
\fancyhf{}
\fancyfoot{}
\renewcommand{\headrulewidth}{0pt}
\renewcommand{\footrulewidth}{0pt}
\setlength{\footskip}{8pt}

\addtolength{\oddsidemargin}{-0.5in}
\addtolength{\evensidemargin}{-0.5in}
\addtolength{\textwidth}{1in}
\addtolength{\topmargin}{-.5in}
\addtolength{\textheight}{1.0in}

\urlstyle{same}
\raggedright
\setstretch{1.15}
\setlength{\parindent}{0pt}
\setlength{\parskip}{11pt}

\begin{document}

\vspace{-10pt}
\begin{tabular*}{\textwidth}{l@{\extracolsep{\fill}}r}
  \textbf{\Large Jan J. Serra} & \small \href{https://janjs.dev}{\underline{janjs.dev}} \\
\end{tabular*}
\vspace{-2pt}

\noindent
{\scriptsize
  {\color{Gray}GitHub\color{black}} \href{https://github.com/janjs}{\underline{github.com/janjs}}
  \hspace{0.35em} $|$ \hspace{0.35em}
  {\color{Gray}Phone\color{black}} +31 683711678
  \hspace{0.35em} $|$ \hspace{0.35em}
  {\color{Gray}Email\color{black}} \href{mailto:jan.jime.serra@gmail.com}{jan.jime.serra@gmail.com}
  \hspace{0.35em} $|$ \hspace{0.35em}
  {\color{Gray}Location\color{black}} Utrecht, Netherlands
}

\vspace{18pt}
\noindent
{\small 20 August 2026}

\vspace{6pt}
\noindent
Dear Remote team,

I'm an AI software engineer with more than five years of experience turning unclear business problems into production systems. I'm applying for the Senior Forward Deployed Engineer role because it describes the work I enjoy most: speaking directly with people who have a real problem, working out what actually needs to happen, and owning the solution through design, build and rollout.

At Accenture, I've worked as a consultant embedded with teams at Merz Pharma, Medtronic and Rabobank. At Merz, I currently own the end-to-end delivery of a production multi-agent platform. I work directly with the client to identify useful automation, turn incomplete requirements into technical designs, and ship agents and enterprise integrations into production. The platform includes RAG, tool calling, streaming chat, workflow orchestration, Microsoft's Copilot ecosystem and Google's Agent2Agent protocol.

I like experimenting with AI, but demos by themselves do not hold my interest for very long. At Medtronic, I helped build an AI platform used by more than 50,000 employees. We had to spread traffic across Azure OpenAI deployments in several regions and make long-running generations survive when a user changed chats. Before AI, I worked on Rabobank identity services that handled more than 1,000 requests per second with 99.99\% availability. That work taught me not to treat reliability as something to add later.

I've always been more comfortable building things than presenting them. Over the past few years, I've pushed myself to change that. One of my first internal talks came soon after I joined Accenture, when GitHub Copilot had just come out. Most people around me did not expect tools like it to become part of everyday coding. More recently, I gave an in-person talk about MCP. I still get nervous before speaking in front of a room, but that is part of why doing it has been so rewarding. It has made me better at explaining technical ideas and more willing to take on things that make me uncomfortable.

Remote's product also makes sense to me personally. I grew up in Barcelona, studied in Dublin and now live in Utrecht. Moving between those countries meant dealing with the paperwork myself.

Consulting has taught me how to enter an unfamiliar company, ask useful questions and earn enough trust to change how people work. What I want now is to stay closer to one product and use what I learn from one customer to make it better for the next. That is what I am looking for.

\vspace{6pt}
\noindent
Best,\\[10pt]
Jan Serra

\end{document}
